% ============================================================================
% COURS 12 — EXERCICES, PARTIE A
% Source : 16_Rdm_2026.docx
% ============================================================================

\section{Exercices du chapitre}
\label{sec:rdm-8}

\begin{center}
\IfFileExists{pandoc/media/image59.png}{\includegraphics[width=.52\linewidth]{pandoc/media/image59.png}}{}
\end{center}

\subsection{Traction -- compression}

\subsubsection{Exercice 1 : choix d'un câble}

\begin{center}
\begin{tabular}{cc}
\IfFileExists{pandoc/media/image61.png}{
\includegraphics[width=.43\linewidth]{pandoc/media/image61.png}}{} &
\IfFileExists{pandoc/media/image62.png}{\includegraphics[width=.43\linewidth]{pandoc/media/image62.png}}{}
\end{tabular}
\end{center}

On souhaite soulever une poutre de béton armé de longueur $6\ \mathrm m$ et de masse $2\ \mathrm t$. Deux sangles sont disposées symétriquement et reliées à une élingue constituée de deux brins de longueur $l=4\ \mathrm m$.

Chaque câble est constitué de six torons de dix-neuf fils de diamètre $0{,}8\ \mathrm{mm}$, entourant une âme textile dont la résistance mécanique est négligée. La limite d'élasticité d'un fil est $R_e=1770\ \mathrm{N/mm^2}$.

\begin{enumerate}
  \item Déterminer la tension maximale dans un brin de câble en fonction de la géométrie de l'élingue.
  \item Le coefficient de sécurité imposé aux structures de levage étant $s=6$, déterminer le diamètre minimal du câble à choisir.
  \item Déterminer le coefficient de sécurité effectif du câble retenu.
\end{enumerate}

\subsubsection{Exercice 2 : câble de puits de mine}

\begin{minipage}{.72\linewidth}
On cherche la forme d'égale résistance d'un câble soumis à son propre poids et supportant à son extrémité une nacelle d'effort $F=2\times10^3\ \mathrm{daN}$.

Données :
\begin{itemize}
  \item poids volumique : $\omega=7{,}8\ \mathrm{daN/dm^3}$ ;
  \item résistance pratique : $R_{pe}=100\ \mathrm{MPa}$ ;
  \item longueur : $L=500\ \mathrm m$.
\end{itemize}
\end{minipage}\hfill
\begin{minipage}{.24\linewidth}
\IfFileExists{pandoc/media/image63.png}{\includegraphics[width=\linewidth]{pandoc/media/image63.png}}{}
\end{minipage}

Pour un câble d'égale résistance, la section varie suivant :
\[
S(x)=\frac{F}{R_{pe}}\exp\!\left(\frac{\omega x}{R_{pe}}\right),
\]
avec des unités cohérentes.

\begin{enumerate}
  \item Déterminer les diamètres aux deux extrémités du câble d'égale résistance.
  \item Déterminer le diamètre d'un câble de section constante supportant la même nacelle.
  \item Calculer et comparer les masses des deux câbles.
\end{enumerate}

\subsubsection{Exercice 3 : bielle de moteur Diesel et chemise mince}

\paragraph{Compression de la bielle.}
Une bielle de moteur Diesel est soumise à un effort maximal de compression de $150\times10^3\ \mathrm N$. Sa partie centrale est modélisée par un prisme de section $S=300\ \mathrm{mm^2}$ et de longueur $L=180\ \mathrm{mm}$.

\begin{enumerate}
  \item Avec un coefficient de sécurité $s=1{,}2$, déterminer la limite d'élasticité minimale de l'acier spécial utilisé.
  \item Le module d'Young étant $E=2\times10^5\ \mathrm{MPa}$, calculer le raccourcissement maximal de la bielle.
\end{enumerate}

\paragraph{Paroi mince de la chemise.}
Le piston a un diamètre intérieur $d=155\ \mathrm{mm}$ et la pression effective dans la chambre vaut $p=8\ \mathrm{MPa}$. La chemise est assimilée à un tube de longueur $L$, de diamètre intérieur $d$ et d'épaisseur $e$ faible devant $d$. La fonte utilisée possède $R_{pe}=50\ \mathrm{MPa}$.

La coupure diamétrale conduit à un effort normal :
\[
N=pLd,
\]
et la section résistante vaut approximativement $2eL$.

\begin{enumerate}[resume]
  \item Déterminer l'épaisseur minimale de la chemise.
\end{enumerate}

\subsubsection{Exercice 4 : maillon de chaîne}

\begin{minipage}{.66\linewidth}
La plaquette d'une chaîne à rouleaux est soumise à la traction. Le perçage du plan $A-A$ produit un coefficient de concentration $K_{tp}=1{,}6$ ; la réduction de section du plan $B-B$ est associée à $K_{tr}=1{,}05$.

Caractéristiques de l'acier :
\[
R_m=1300\ \mathrm{MPa},\qquad
R_e=1100\ \mathrm{MPa},\qquad
E=2{,}2\times10^5\ \mathrm{MPa}.
\]
\end{minipage}\hfill
\begin{minipage}{.30\linewidth}
\IfFileExists{pandoc/media/image65.png}{\includegraphics[width=\linewidth]{pandoc/media/image65.png}}{}
\end{minipage}

\begin{enumerate}
  \item Évaluer la charge limite de rupture statique $F_{\mathrm{lim}}$ en recherchant la section la plus critique.
  \item Pour un fonctionnement dynamique avec $s=3$, déterminer la charge maximale de service $F_{\max}$.
\end{enumerate}

\subsection{Fardeleuse}

\textit{Source : concours commun PT-SI 2011.}

\begin{center}
\begin{tabular}{cc}
\IfFileExists{pandoc/media/image66.jpeg}{\includegraphics[width=.36\linewidth]{pandoc/media/image66.jpeg}}{} &
\IfFileExists{pandoc/media/image67.png}{
\includegraphics[width=.52\linewidth]{pandoc/media/image67.png}}{}
\end{tabular}
\end{center}

Le fardelage consiste à déposer un film plastique thermo-rétractable autour d'un produit puis à le chauffer. Le produit est déplacé par un tapis entraîné par adhérence entre un cylindre moteur et un cylindre tendeur.

La tension initiale $T_0$ est réglée en translatant l'axe du cylindre tendeur. Le cylindre est guidé par deux roulements qui exercent, en $B$ et $C$, deux efforts :
\[
\vec F=-F\vec y,
\qquad F=230\ \mathrm N.
\]
L'axe est modélisé par une poutre cylindrique de diamètre $d$, reposant sur deux appuis en $A$ et $D$.

\begin{center}
\begin{tabular}{cc}
\IfFileExists{pandoc/media/image68.png}{\includegraphics[width=.38\linewidth]{pandoc/media/image68.png}}{} &
\IfFileExists{pandoc/media/image69.png}{\includegraphics[width=.50\linewidth]{pandoc/media/image69.png}}{}
\end{tabular}
\end{center}

\subsubsection{Dimensionnement de l'axe}

\begin{enumerate}
  \item Établir la relation entre la tension $T_0$ du tapis et la norme $F$ des efforts en $B$ et $C$.
  \item Déterminer les réactions des appuis en $A$ et $D$ en fonction de $F$.
  \item Exprimer les composantes du torseur de cohésion le long de la poutre et tracer les diagrammes correspondants. Identifier la sollicitation.
  \item En négligeant l'effort tranchant, localiser la zone la plus sollicitée.
  \item Déterminer la contrainte normale maximale $\sigma_{xx,\max}$.
  \item Avec un coefficient de sécurité $s_c=2$, établir l'expression du diamètre minimal $d_{\min}$.
\end{enumerate}

\subsubsection{Choix du matériau}

Les barres disponibles ont un diamètre compris entre $3$ et $50\ \mathrm{mm}$.

\begin{center}
\small
\begin{tabular}{lllrr}
\toprule
Acier & État & & $R_e$ (MPa) & $R_m$ (MPa)\\
\midrule
C22 & normalisé && 240 & 430\\
C22 & trempé et revenu && 340 & 650\\
C55 & normalisé && 370 & 680\\
C55 & trempé et revenu && 550 & 950\\
15 Cr Ni 6 & trempé et revenu && 650 & 1000\\
25 Cr Mo 4 & trempé et revenu && 550 & 850\\
\bottomrule
\end{tabular}
\end{center}

\begin{enumerate}[resume]
  \item Choisir un acier et un état métallurgique compatibles avec le diamètre disponible et le coefficient de sécurité.
  \item Pour l'acier choisi, calculer $d_{\min}$ puis sélectionner le diamètre commercial immédiatement supérieur.
\end{enumerate}
